%% chapters/ch3_method.tex

\chapter{研究方法}
\label{ch:method}

\section{系統架構}
\label{sec:architecture}

這一節說明整體流程。並排圖請放在同一個 \verb|figure| 環境裡，
用 \verb|\mpfig| 系列巨集，中間以 \verb|\hfill| 分隔：

\begin{figure}[htbp]
  \centering
  \mpfig{0.45}{figures/example_fig}{左側子圖說明。}{fig:panel-left}%
  \hfill
  \mpfig{0.45}{figures/example_fig}{右側子圖說明。}{fig:panel-right}
\end{figure}

圖 \ref{fig:panel-left} 與圖 \ref{fig:panel-right} 並排顯示。

\section{實驗設計}
\label{sec:experiment-design}

表格請用一般的 \verb|table| 環境，表的標題一律放在表格上方：

\begin{table}[htbp]
  \centering
  \caption{示範用表格。實際撰寫時請填入自己的資料。}
  \label{tab:example}
  \begin{tabular}{lrr}
    \hline
    項目 & 數值一 & 數值二 \\
    \hline
    條件 A & 0.00 & 0.00 \\
    條件 B & 0.00 & 0.00 \\
    \hline
  \end{tabular}
\end{table}

表 \ref{tab:example} 列出實驗條件。
